% TOP_PROBLEM: {"id": 2900, "title": "Kirby Problem 4.24", "category_id": 7, "category": {"id": 7, "name": "topology", "display_name": "Topology", "description": "Properties preserved under continuous deformations.", "slug": "topology", "order_index": 7, "created_at": "2026-07-31T15:26:25.671Z"}, "status": "open"}
% FIRST_POSTED: null
% ENTRY_KIND: statement_only
% Imported from: batch13.tex; UnsolvedMath ID 2900
\documentclass[11pt]{article}
\usepackage[margin=25mm]{geometry}
\usepackage{fontspec}
\setmainfont{Latin Modern Roman}
\usepackage{amsmath,amssymb,mathtools}
\usepackage{unicode-math}
\setmathfont{Latin Modern Math}
\usepackage{xurl,hyperref}
\hypersetup{colorlinks=true,urlcolor=blue}
\setcounter{secnumdepth}{0}
\setlength{\parindent}{0pt}
\setlength{\parskip}{5pt}
\setlength{\emergencystretch}{3em}
\newcommand{\sref}[2]{\hyperlink{source:#1:#2}{[S#2]}}
\newcommand{\cataloguescope}[1]{\paragraph{Recorded target.}#1}
\begin{document}
% UnsolvedMath ID 2900; source catalogue code KP-4.24
\setcounter{section}{2899}
\section{Kirby Problem 4.24}\label{2900}
{\small\sffamily UnsolvedMath ID 2900 · Topology}
\subsection{Definitions and mathematical statement}

\paragraph{Shared notation.}$\mathbb N=\{1,2,\ldots\}$, and $\mathbb Z,\mathbb Q,\mathbb R,\mathbb C$ denote the integers, rationals, reals and complex numbers. Any other conventions are specified locally.


Let $K\subset S^3=\partial B^4$ be a smooth knot with an integer framing $n$. Let $W(K,n)$ be its two-handle trace
\[
W(K,n)=B^4\cup_{\nu K}(D^2\times D^2),
\]
where the attaching region is $\partial D^2\times D^2$ and the framing is measured relative to the Seifert longitude of $K$. The cocore is the disk $C=\{0\}\times D^2$ in the attached handle. Choose a meridian $U$ of $K$ in the attaching-boundary diagram and regard it as a curve in $\partial W(K,n)$. A carving disk is a smooth proper embedded disk $D\subset W(K,n)$ with $\partial D=U$. Arrange the interiors of $D$ and $C$ to be transverse and their boundaries disjoint. Their algebraic intersection is the signed count
\[
D\cdot C=\sum_{p\in\operatorname{int}D\cap\operatorname{int}C}\epsilon_p,
\qquad\epsilon_p\in\{1,-1\},
\]
using chosen disk orientations and the orientation of the trace. Its vanishing does not require geometrically disjoint interiors. Remove an open normal disk-bundle neighborhood of $D$ and smooth corners; denote the resulting manifold by $B_D$.
\paragraph{Statement recorded in the supplied source.}
Can this carving construction produce a nonstandard four-ball? In the framed-knot setting above, is there a carving disk $D$ such that
\[
D\cdot C=0\qquad\text{and}\qquad B_D\not\cong_{\mathrm{diff}}B^4?
\]
The disk may enter the two-handle; it is not required to lie in the original zero-handle. The intersection condition is the geometric signed-count condition with the specified cocore. \sref{2900}{2}
\subsection{Short English statement}
Can carving a disk from a framed knot trace give a manifold not diffeomorphic to the four-ball when the disk has zero algebraic intersection with the handle cocore?
\subsection{Sources}
\begin{description}
\item[\hypertarget{source:2900:1}{S1}] ULAM.AI and the UnsolvedMath Contributors, UnsolvedMath: A Curated Collection of Open Mathematics Problems, record 2900 (KP-4.24). CC BY 4.0. \url{https://huggingface.co/datasets/ulamai/UnsolvedMath}.
\item[\hypertarget{source:2900:2}{S2}] R. İnanç Baykur, Robion C. Kirby and Daniel Ruberman, K3—A New Problem List in Low-Dimensional Topology, Mathematical Surveys and Monographs 295 (2026), author version, Problem 4.24 and Remarks, PDF page 210. \url{https://math.berkeley.edu/sites/default/files/surv-295-ruberman-watermarked-author-pdf.pdf}.
\end{description}
\label{end:2900}
\end{document}
